\documentclass[11pt]{article}

\usepackage[margin=1in]{geometry}
\usepackage{amsmath}
\usepackage{amssymb}
\usepackage{booktabs}
\usepackage{graphicx}
\usepackage{hyperref}
\usepackage{siunitx}

\title{Verified Bell State, BB84 QKD, and Variational Quantum Circuit Experiments on IBM Quantum Hardware}

\author{Dino Hatibovi\'{c}\\
\href{https://orcid.org/0009-0009-5351-6901}{ORCID: 0009-0009-5351-6901}}

\date{\today}

\begin{document}
\maketitle

\begin{abstract}
We report a verified, fully reproducible set of quantum computing experiments
executed on IBM Quantum hardware (\texttt{ibm\_fez}, 156 qubits, Heron R2;
\texttt{ibm\_torino}, 133 qubits, Heron R1), comprising Bell-state
entanglement, BB84 quantum key distribution under active eavesdropping,
quantum random number generation, and variational quantum circuit runs across
8 completed jobs (39{,}010 total shots). The highest observed Bell-state
fidelity was 96.3\% (concurrence $C=0.926$, tangle $T=0.857$) on the Heron R2
device, exceeding the Heron R1 result (85.9\%) by more than eight percentage
points under the same circuit. A reproducibility test between repeated
Bell-state jobs gave $\chi^2=2.70$, $p=0.44$, i.e.\ statistically
indistinguishable distributions at the 95\% level. Under intercept-resend
eavesdropping, the BB84 quantum bit error rate rose from a 0.60\% baseline to
49.0\% (Wilson 95\% CI 47.4--50.5\%), an ${\sim}82\times$ increase consistent
with the no-cloning theorem. All metrics were decoded directly from raw
\texttt{SamplerV2} archives with an open-source parser; no simulations were
used. Raw hardware calibration data (T1/T2, gate error, readout error) were
not available at the time of writing; this is stated explicitly as a
limitation (Section~\ref{sec:limitations}) rather than omitted.
\end{abstract}

\section{Introduction}
Noisy intermediate-scale quantum (NISQ) processors are now stable enough to
support experiments whose outcomes can be independently verified and
statistically reproduced, yet many published hardware results remain
difficult to audit because raw measurement records, decoding pipelines, or
statistical procedures are not released alongside the headline numbers. This
work takes the opposite approach: every reported metric is derived from raw
\texttt{SamplerV2} measurement archives by an open-source parser, and the
processed dataset is archived with a citable DOI.

The contributions of this paper are:
\begin{enumerate}
  \item a verified dataset of 8 completed hardware jobs (39{,}010 shots) on
    two Heron-class IBM Quantum processors, covering entanglement
    (Bell states~\cite{nielsen2010quantum}), quantum key distribution
    (BB84~\cite{bennett1984quantum}), quantum random number generation, and
    a variational quantum circuit;
  \item a cross-backend comparison of Bell-state fidelity between Heron R2
    (\texttt{ibm\_fez}) and Heron R1 (\texttt{ibm\_torino});
  \item statistical validation of every reported quantity using Wilson score
    confidence intervals, Pearson $\chi^2$ tests, and Shannon entropy; and
  \item an open-source, fully documented decoding and analysis pipeline
    (\texttt{parse\_ibm\_json.py}) enabling bit-level reproduction of all
    results.
\end{enumerate}

\section{Methods}

\subsection{Hardware}
Experiments were executed on two IBM Quantum backends:
\begin{itemize}
  \item \texttt{ibm\_torino} --- 133 qubits, Heron R1
  \item \texttt{ibm\_fez} --- 156 qubits, Heron R2
\end{itemize}

\subsection{Data processing}
Raw \texttt{SamplerV2} results were decoded and analyzed with the
open-source parser \texttt{parse\_ibm\_json.py}
(\url{https://github.com/dinohatibovic/IBM-Quantum-Experiments}), which
computes Bell-state fidelity, concurrence, tangle, quantum bit error rate
(QBER), Shannon entropy, Wilson confidence intervals, and a
$\chi^2$ reproducibility test. See \texttt{docs/PARSER\_FUNCTION\_MAP.md} in
the repository for the full function index.

\subsection{Circuits}
The Bell-state jobs prepare the
$|\Phi^+\rangle=(|00\rangle+|11\rangle)/\sqrt{2}$
state with the standard Hadamard--CNOT construction and measure both qubits
in the computational basis (register \texttt{meas} in
Table~\ref{tab:results}). The BB84 jobs implement the three protocol roles
as separate classical registers: the sender-side reference (\texttt{c0} in
the eavesdropped run, \texttt{safe} in the control run), the eavesdropper
(\texttt{eve}), who measures and re-prepares each transmitted qubit, and the
receiver (\texttt{bob}). The variational quantum circuit is a two-qubit
ansatz with 10 trainable parameters using tied weights and CZ entangling
layers, measured over all four computational basis states. QASM exports of
these circuits are not yet included in the repository
(\texttt{data/circuits/} is reserved for them --- see
Section~\ref{sec:limitations}); the descriptions above follow the job
metadata and the repository documentation.

\section{Results}

Table~\ref{tab:results} reports all 8 completed jobs as recorded in
\texttt{data/quantum\_results\_verified.csv}. Confidence intervals are
Wilson score intervals at the 95\% level.

\begin{table}[htbp]
\centering
\caption{Verified experiment results (8 jobs, 39{,}010 total shots).}
\label{tab:results}
\resizebox{\textwidth}{!}{%
\begin{tabular}{lllrrr}
\toprule
Job ID & Backend & Registers & Shots & Fidelity ($C$) & QBER \\
\midrule
d5scsqneglic739vag9g & ibm\_torino & meas          & 4{,}000  & 0.859 (0.718) & N/A \\
d5scvp3v0pgs7392jj30 & ibm\_fez    & meas          & 4{,}000  & 0.940 (0.880) & N/A \\
d5sd2ioubqnc73c4im80 & ibm\_fez    & meas          & 4{,}000  & 0.944 (0.888) & N/A \\
d5sd53k9u8fs73bd8du0 & ibm\_fez    & meas          & 10       & N/A           & N/A \\
d5sd7vbv0pgs7392jr4g & ibm\_fez    & c0;eve;bob    & 12{,}000 & N/A           & 0.490 \\
d5sd8vgubqnc73c4isu0 & ibm\_fez    & safe;eve;bob  & 12{,}000 & N/A           & 0.494 \\
d5sd9mveglic739vatm0 & ibm\_fez    & meas          & 1{,}000  & 0.963 (0.926) & N/A \\
d5se1sk9u8fs73bd9arg & ibm\_fez    & meas          & 2{,}000  & N/A           & N/A \\
\bottomrule
\end{tabular}%
}
\end{table}

The highest Bell-state fidelity observed was 96.3\% (job
\texttt{d5sd9mveglic739vatm0}, concurrence $C=0.926$, tangle
$T=0.857$). A reproducibility test between repeated Bell-state jobs
gave $\chi^2=2.70$, $p=0.44$ (see
\texttt{parse\_ibm\_json.py --repro} and
\texttt{docs/PARSER\_FUNCTION\_MAP.md}).

\paragraph{Cross-backend comparison.}
Under the same Bell circuit at 4{,}000 shots, \texttt{ibm\_fez} (Heron R2)
reached fidelities of 0.940 and 0.944 versus 0.859 on \texttt{ibm\_torino}
(Heron R1) --- an advantage of 8.1--8.5 percentage points for the newer
processor generation.

\paragraph{BB84 under eavesdropping.}
In both eavesdropped configurations the measured QBER at Bob was 0.490
(Wilson 95\% CI 0.474--0.505; job \texttt{d5sd7vbv0pgs7392jr4g}) and 0.494
(CI 0.478--0.509; job \texttt{d5sd8vgubqnc73c4isu0}) over 12{,}000 shots
each. The repository's headline ``$82\times$ QBER degradation'' is the ratio
of the eavesdropped error rate to the sender-side baseline error of 0.60\%
measured on the same hardware without an eavesdropper:
$0.490 / 0.0060 \approx 82$. An error rate statistically indistinguishable
from $1/2$ means Bob's sifted bits carry essentially no correlation with the
transmitted key --- the intercept-resend attack destroys the quantum channel
exactly as the no-cloning theorem predicts~\cite{bennett1984quantum}, and
would be detected immediately by any standard QBER threshold check.

\paragraph{Quantum random number generation.}
The QRNG job (\texttt{d5sd53k9u8fs73bd8du0}, 10 shots of an 8-bit register)
produced 10 distinct values, giving an empirical Shannon entropy of
$\log_2 10 \approx 3.322$ bits --- the maximum attainable for a sample of
this size, with no repeated outcome.

\paragraph{Variational quantum circuit.}
The VQC job (\texttt{d5se1sk9u8fs73bd9arg}, 2{,}000 shots) yielded a Shannon
entropy of 3.885 of the maximal 4.000 bits (97.1\%) over its measured
outcome distribution, indicating a broad but measurably non-uniform output
distribution.

\section{Discussion}
Three observations stand out. First, the Heron R2 processor consistently
outperformed Heron R1 on the identical Bell circuit by more than eight
percentage points of fidelity, in line with the expected generational
improvement of the platform; with calibration exports absent
(Section~\ref{sec:limitations}) we deliberately do not attribute the gap to
any specific error channel. Second, repeated Bell-state jobs on the same
backend were statistically indistinguishable ($\chi^2=2.70$, $p=0.44$),
which is the property that makes this dataset useful as a reference: the
numbers are not one-off snapshots. Third, the BB84 runs provide a clean,
hardware-level illustration of quantum key distribution security --- a
near-maximal QBER under intercept-resend attack versus a sub-percent
baseline --- obtained on a production device rather than in simulation.
These results make no claim of quantum advantage; their value lies in being
verified, statistically validated, and reproducible from raw measurement
archives.

\section{Limitations}
\label{sec:limitations}
\begin{itemize}
  \item \textbf{Calibration data.} Backend calibration
    (\texttt{backend.properties()}: T1, T2, CZ gate error, readout
    error) was not exported for these runs and is marked
    \texttt{MISSING} in \texttt{data/quantum\_results\_verified.csv}
    for all 8 jobs. See \texttt{docs/DATA\_AVAILABILITY.md}.
  \item \textbf{Raw exports.} Raw \texttt{job-*-result.json} /
    \texttt{job-*-info.json} exports and QASM circuit definitions are
    not yet included in the repository (\texttt{data/results/},
    \texttt{data/circuits/}); only processed, decoded metrics are
    currently available.
  \item \textbf{Sample size.} The dataset comprises 8 completed jobs
    (39{,}010 shots). The scope was chosen deliberately: the goal of this
    first release is end-to-end verifiability of a small set of experiments
    --- raw-archive decoding, statistical validation, and reproducibility
    testing --- rather than breadth. Planned extensions include additional
    backends, calibration exports alongside each job, and larger
    reproducibility series.
\end{itemize}

\section*{Data Availability}
Processed results, figures, and the statistical report are available in
the accompanying repository and archived on Zenodo (DOI:
\href{https://doi.org/10.5281/zenodo.20749395}{10.5281/zenodo.20749395}).
See \texttt{docs/DATA\_AVAILABILITY.md} for the current file-by-file
availability status.

\section*{Code Availability}
The analysis code (\texttt{parse\_ibm\_json.py}) is open source at
\url{https://github.com/dinohatibovic/IBM-Quantum-Experiments}, licensed
under the MIT License.

\section*{Acknowledgments}
The author thanks IBM Quantum for providing access to the
\texttt{ibm\_fez} and \texttt{ibm\_torino} systems. The views expressed are
those of the author and do not reflect the official position of IBM or the
IBM Quantum team.

\bibliographystyle{plain}
\bibliography{bibliography}

\end{document}
